% Part: computability
% Chapter: computability-theory
% Section: normal-form

\documentclass[../../../include/open-logic-section]{subfiles}

\begin{document}

\olfileid{cmp}{thy}{nfm}
\olsection{স্বাভাবিক রূপের উপপাদ্য}

ধরা যাক, গণনসাধ্য অপেক্ষকের সংজ্ঞা বর্ণনা করা যায় এবং কোনো নিবেশে সেই
অপেক্ষকের গণনার পূর্ণ নথির সম্ভাব্য কোনো বর্ণনা সঠিক কি না পরীক্ষা করা
যায়। তাহলে যুক্তিসঙ্গতভাবেই বলা যায় যে গণনা-মডেল নির্বিশেষে যেকোনো
নিবেশ~$x$-এ যেকোনো গণনসাধ্য অপেক্ষক~$f$-এর মান নিচের উপায়ে নির্ণয় করা
যায়:
\begin{enumerate}
  \item গণনার নথির সম্ভাব্য সব বর্ণনার মধ্যে অনুসন্ধান করা।
  \item কোনো নির্দিষ্ট নথি~$f(x)$-এর গণনার নথি কি না পরীক্ষা করা।
  \item সঠিক নথিটি পাওয়া গেলে তা থেকে~$f(x)$-এর মান বার করা।
\end{enumerate}
বাস্তবে যে ঠিক এমনটিই হয়, সেটিই ক্লিনির স্বাভাবিক রূপের উপপাদ্যের
বক্তব্য।

\begin{thm}[ক্লিনির স্বাভাবিক রূপের উপপাদ্য]
\ollabel{thm:normal-form}
এমন একটি আদিম পুনরাবৃত্ত সম্বন্ধ~$T(e, x, s)$ এবং একটি আদিম পুনরাবৃত্ত
অপেক্ষক~$U(s)$ আছে, যাদের নিচের ধর্ম রয়েছে: $f$ যেকোনো আংশিক গণনসাধ্য
অপেক্ষক হলে কোনো~$e$-এর জন্য
\[
f(x) \simeq U(\umin{s}{T(e, x, s)})
\]
প্রত্যেক~$x$-এর ক্ষেত্রে সত্য।
\end{thm}

\begin{proof}[প্রমাণের রূপরেখা]
যেকোনো গণনা-মডেলের জন্য গণনসাধ্য অপেক্ষক~$f$-এর একটি বর্ণনা বিধিবদ্ধভাবে
সংজ্ঞায়িত করা যায় এবং সেই বর্ণনাকে কোনো স্বাভাবিক সংখ্যা~$e$ দিয়ে
সংকেতায়িত করা যায়। একইভাবে, সূচক~$e$-এর অপেক্ষকটির নিবেশ~$x$-এ গণনার
প্রক্রিয়াটি লিপিবদ্ধ করে এমন “গণনা-অনুক্রম”-এর ধারণা বিধিবদ্ধভাবে
সংজ্ঞায়িত করা যায়। এমন গণনা-অনুক্রমকেও কোনো সংখ্যা~$s$ দিয়ে সংকেতায়িত
করা যায়। কাজটি এমনভাবে করা যায় যাতে
\begin{enumerate}
  \item সম্বন্ধ $T(e, x, s)$ তখনই সত্য হয়, যখন সংখ্যা~$s$ নিবেশ~$x$-এ
  সূচক~$e$-এর অপেক্ষকটির গণনা-অনুক্রমকে সংকেতায়িত করে; এবং
  \item অপেক্ষক $U(s)$ সংকেত~$s$-এর গণনা-অনুক্রমকে সেই গণনার চূড়ান্ত
  ফলে পাঠায়—
\end{enumerate}
দুটিই গণনসাধ্য হয়। আসলে সম্বন্ধ~$T$ ও অপেক্ষক~$U$ উভয়ই আদিম পুনরাবৃত্ত।
\end{proof}

\begin{explain}
স্বাভাবিক রূপের উপপাদ্যের বিধিবদ্ধ প্রমাণ দিতে হলে কোনো একটি গণনা-মডেল
স্থির করতে হবে, গণনসাধ্য অপেক্ষকের বর্ণনা ও গণনা-অনুক্রমের সংকেতায়ন
বিশদে সম্পন্ন করতে হবে এবং যাচাই করতে হবে যে সম্বন্ধ~$T$ ও অপেক্ষক~$U$
আদিম পুনরাবৃত্ত। অধিকাংশ প্রয়োগের জন্য $T$ ও~$U$ গণনসাধ্য এবং~$U$
সর্বত্র-সংজ্ঞায়িত—এইটুকুই যথেষ্ট।

স্বাভাবিক রূপের উপপাদ্যের প্রমাণটি এই সংক্ষিপ্ত বাক্যে মনে রাখাই সম্ভবত
সবচেয়ে ভালো: $\umin{s}{T(e, x, s)}$ নিবেশ~$x$-এ সূচক~$e$-এর অপেক্ষকটির
একটি গণনা-অনুক্রম খোঁজে, আর সেটি পাওয়া গেলে $U$ গণনা-অনুক্রমটির নির্গম
ফেরত দেয়।
\end{explain}

গণনা-মডেলটি আংশিক পুনরাবৃত্ত অপেক্ষক হলে (ক্লিনি মূলত এটিই ব্যবহার
করেছিলেন), উপপাদ্যটি দেখায় যে যেকোনো অপেক্ষক সংজ্ঞায়িত করতে সীমাহীন
অনুসন্ধান মাত্র একবার, অর্থাৎ একটি মাত্র $\umin{y}{f(\vec x, y)}$
অপারেটর ব্যবহার করাই যথেষ্ট। এই অর্থে প্রত্যেক আংশিক পুনরাবৃত্ত অপেক্ষকের
একটি স্বাভাবিক রূপ আছে।

$T$ ও $U$ ব্যবহার করে $\cfind{0}$, $\cfind{1}$, $\cfind{2}$, \dots
তালিকাটি সংজ্ঞায়িত করা যায়। এখন থেকে আমরা ধরে নেব, $T$ ও~$U$-এর একটি
উপযুক্ত নির্বাচন স্থির করা হয়েছে, এবং সমীকরণ
\[
\cfind{e}(x) \simeq U(\umin{s}{T(e,x,s)})
\]
হলো~$\cfind{e}$-এর \emph{সংজ্ঞা}।

আরেকটি উপযোগী তথ্য হলো:

\begin{thm}
প্রত্যেক আংশিক গণনসাধ্য অপেক্ষকের অসীমসংখ্যক সূচক আছে।
\end{thm}

এটিও স্বজ্ঞাতভাবে স্পষ্ট। কোনো গণনসাধ্য অপেক্ষকের (বর্ণনা) দেওয়া থাকলে
একই অপেক্ষক (একই নিবেশ-নির্গম জোড়া) গণনা করে এমন ভিন্ন একটি বর্ণনা
তৈরি করা যায়—যেমন, প্রথমে গণনায় কোনো প্রভাব ফেলে না এমন কিছু করা
(ধরা যাক, $0 = 0$ কি না পরীক্ষা করা বা $5$ পর্যন্ত গোনা), তারপর মূল
গণনাটি করা। পরিবর্তিত বর্ণনার সূচকটি সবসময় মূল সূচক থেকে আলাদা হবে।
দুটিই একই অপেক্ষকের সূচক; শুধু অপেক্ষকটিকে সামান্য ভিন্নভাবে গণনা করে।

\end{document}
